\begin{textsample}{POS Dim 8 – human – Score 2.00 – es/es\_ef\_anos\_iniciais\_arte\_5\_p10.txt}  \label{ex:f8_pos_016}
Nos quadros a seguir, que apresentam as unidades temáticas, os objetos de conhecimento e as habilidades definidas para cada ano ( ou \textbf{bloco} de anos ), cada habilidade é identificada por um código alfanumérico cuja composição é a \textbf{seguinte} Vale destacar que o uso de numeração sequencial para identificar as habilidades de cada ano ou \textbf{bloco} de anos não representa uma ordem ou hierarquia esperada das aprendizagens ( BRASIL, 2017 ).
Também é importante ressaltar que as habilidades representam o que se espera que os estudantes aprendam ao longo de cada ano do Ensino Fundamental e não descrevem ações ou condutas docentes, nem induzem à opção por abordagens ou metodologias, que devem ser adotadas, adequando -se à \textbf{realidade} de cada unidade de ensino, considerando o contexto e as características de seus estudantes ( BRASIL, 2017 ).

% matched lemmas: bloco, realidade, seguinte
\end{textsample}
